\documentclass[a4paper]{article}

\usepackage{mathrsfs}
% ============================================
%         General Packages
% ============================================
\usepackage{balance}

\usepackage{kotex}      % for Korean
\usepackage{lipsum}     % for dummy text
\usepackage{array}

\usepackage{cite}   % for citation

\usepackage{amsmath,amsfonts}
\usepackage{amssymb}
\usepackage{soul}
\usepackage{mathrsfs}

\usepackage{hyperref}

    \hypersetup{
        pdftoolbar=false,        	% show Acrobat’s toolbar?
        pdfmenubar=false,        	% show Acrobat’s menu?
        pdffitwindow=false,     	% window fit to page when opened
        pdfstartview={FitH},    	% fits the width of the page to the window
        pdftitle={},    	% title
        pdfauthor={},     	% author
        %     pdfsubject={},   	% subject of the document
        %     pdfcreator={},   	% creator of the document
        %     pdfproducer={}, 	% producer of the document
        %     pdfkeywords={keyword1, key2, key3}, % list of keywords
        %     pdfnewwindow=true,      	% links in new PDF window
        colorlinks=true,       	% false: boxed links; true: colored links
        linkcolor=blue,          	% color of internal links (change box color with linkbordercolor)
        citecolor=blue,       	% color of links to bibliography
        filecolor=blue,      	% color of file links
        urlcolor=blue		% color of external links
    }

\usepackage{algorithm}
\usepackage{algorithmic}

\usepackage{epsfig}
\usepackage[dvipsnames]{xcolor}

% MY PACKAGES ===================================

\usepackage{svg}
\usepackage{subfig}

\usepackage{textcomp}
\usepackage{stfloats}
\usepackage{url}
\usepackage{verbatim}
\usepackage{graphicx}
\usepackage{multirow}
\usepackage{multicol}

\let\proof\relax
\let\endproof\relax
\usepackage{amsthm}
\usepackage{lipsum}
\usepackage{tikz}

% 
\usepackage{titlepic}
\usepackage{pdfpages}
\usepackage{etoolbox} % 조건문을 쓰기 위한 패키지


\titlepic{%
  
\includegraphics[height = 8mm]{figures/logo_KAIST_simple.png}
  \hspace{5mm}
  
\includegraphics[height = 8mm]{figures/MIC_Logo.png}
  }

% ========================================
%         TITLE and AUTHORS
% ========================================
\title{
  Real-World RL
}

\author{
    Hyochan Lee
    \thanks{\href{https://kaist-mic-lab.github.io/members/hcLee/}{Hyochan Lee} is Ph.D. student at KAIST, {\tt\small hyochanlee@kaist.ac.kr}}%
    , and
    Kyunghwan Choi
    \thanks{\href{https://kaist-mic-lab.github.io/members/khChoi/}{Kyunghwan Choi} is Assistant Professor at KAIST}%
}

\date{
    Compiled on \today
}

\begin{document}

\maketitle

\begin{figure}[h]
  \centering
  \includegraphics[width=\textwidth]{figures/01_real_world_rl.png}
  \caption{Real interactions generate streaming transitions that train both the critic and policy online
.}
  \label{fig:}
\end{figure}


\nocite{*}

The environment returns measured transition records; an assumed state
estimator converts the measurements into $\widehat x_k$ before the
state-based transition enters $\mathcal S_t$. Without an estimator,
$\widehat x_k$ is replaced by an information state $o_k$. Physical interaction
is indexed by $k$, whereas network updates are indexed by $t$; the two clocks
need not advance at the same rate. Here $k(t)$ is the latest physical-state
sample available at learner update $t$, and $t(k)$ is the latest learner
update available at physical step $k$.

\section{Problem Definition}

\href{https://kaist-mic-lab.github.io/research/research_themes/00_online_learning_based_optimal_control/}{Online Learning-Based Optimal Control}
defines the desired policy through Bellman optimality. The broader Real-World
RL Theme addresses the case in which its critic and, when explicitly
parameterized, policy are learned from real closed-loop data. A representative
approximate policy is
\begin{equation}
  \widehat\mu_{\phi_t}(\widehat x)
  \approx
  \arg\min_{u\in\mathcal U(\widehat x)}
  \left\{
  g(\widehat x,u)
  +
  \alpha\widehat J_{\theta_t}
  \!\left(\widehat f_{\psi_t}(\widehat x,u)\right)
  \right\}.
\end{equation}

Here $\widehat J_{\theta_t}$ is the critic, $\widehat\mu_{\phi_t}$ is the
policy network, $\widehat f_{\psi_t}$ is a successor model with parameter
$\psi_t$, and $0<\alpha<1$ is the discount factor. The expression above is
model-based. Without a successor model, an observed transition provides a
sample-based Q-factor target directly:
\begin{equation}
  \widehat Q_t^{\mathrm{tar}}(\widehat x_k,u_k)
  :=
  g(\widehat x_k,u_k)
  +
  \alpha\widehat J_{\theta_t}(\widehat x_{k+1}).
\end{equation}
A learned Q-critic fits such targets over state-action pairs; model-free policy
improvement then minimizes the learned Q-factor over $u$.

For the primary notation, a state estimator is assumed to supply the
controller-facing state estimate. The real system then supplies a stream
\begin{equation}
  d_k=(\widehat x_k,u_k,g_k,\widehat x_{k+1}),
  \qquad
  \mathcal D_{0:k(t)-1}
  :=
  (d_0,\ldots,d_{k(t)-1}),
  \qquad
  \mathcal S_t\subseteq\mathcal D_{0:k(t)-1}.
\end{equation}
Here $g_k:=g(\widehat x_k,u_k)$ is the realized stage cost. Because a complete
transition ending at the latest state sample $k(t)$ is indexed by $k(t)-1$,
$\mathcal D_{0:k(t)-1}$ is the ordered collection of all completed transition
records available to the learner. The update set $\mathcal S_t$ may contain
the latest transition alone or selected prior transitions. Replay or minibatch
training is optional rather than a defining assumption.

Without a state estimator, the problem becomes partially observed. Then
$\widehat x_k$ should be replaced by an observation or information state
$o_k$: $o_k=y_k$ only when the current measurement is sufficient for control;
otherwise $o_k$ should summarize an observation history, belief state, or
learned latent state. The corresponding record is
$d_k=(o_k,u_k,g_k,o_{k+1})$.

In an actor-critic realization, the selected update data train both learned
objects:
\begin{equation}
  \left(
  \widehat J_{\theta_t},
  \widehat\mu_{\phi_t}
  \right)
  \xrightarrow[\mathcal S_t]{\mathrm{online\ training}}
  \left(
  \widehat J_{\theta_{t+1}},
  \widehat\mu_{\phi_{t+1}}
  \right),
  \qquad
  u_k=\widehat\mu_{\phi_{t(k)}}(\widehat x_k).
\end{equation}

\section{Why This Problem Matters}

Many control-oriented RL methods are trained and evaluated primarily in
simulation. Their deployed performance can then be limited by the sim-to-real
gap: unmodeled dynamics, sensing imperfections, disturbances, and operating
conditions absent from the simulator.

Learning directly from the real system can capture this missing information
and continue improving when deployment reveals new operating conditions. Here,
adaptability means accumulating and reusing experience to learn a critic and
policy that remain useful over a declared task domain. It does not mean only
tracking a local parameter change along the currently excited trajectory, as
in an adaptive-control-like update.

\section{Key Challenges}

\begin{itemize}
  \item \textbf{Limited and policy-dependent data:} real transitions arrive
  sequentially and are correlated. Unlike simulator rollouts, they cannot be
  generated freely over the entire state-action domain.
  \item \textbf{Local fit versus task-level learning:} reducing loss on the
  latest trajectory can produce a local fit, forget earlier regimes, or
  repeatedly relearn the same task. The learned functions must retain broader
  task-relevant structure.
  \item \textbf{Learning inside the closed loop:} critic and policy change
  while the system is operating. Update time, action latency, stability,
  constraints, fallback behavior, and safety evidence must therefore be part
  of the learning problem.
\end{itemize}

\section{A Conventional Approach: TD-Error-Based Approximate Policy Iteration}

One conventional approach is TD-error-based approximate policy iteration: fix
the current policy, reduce its one-step temporal-difference error, and then
perform policy improvement. For notational simplicity, this section writes
$x_k$ as the state available to the controller; when an estimator is used, it
is replaced by $\widehat x_k$. Here $\pi_i$ is the policy at policy-iteration
step $i$.
\begin{equation}
  \begin{aligned}
    \delta_k^{\pi_i}
    &=
    g(x_k,\pi_i(x_k))
    +
    \alpha\widehat J^{\pi_i}(x_{k+1})
    -
    \widehat J^{\pi_i}(x_k),\\
    \widehat J^{\pi_i}
    &\xleftarrow{\mathrm{TD\ error\ reduction}}
    \operatorname{Fit}\!\left(\delta_k^{\pi_i}\right),
    \qquad
    \pi_{i+1}
    \leftarrow
    \operatorname{Greedy}\!\left(\widehat J^{\pi_i}\right).
  \end{aligned}
\end{equation}

This is a valid approximate-DP mechanism, but streaming data may make it
adaptation-like: a small TD error on visited samples establishes sampled policy
consistency, not task-wide Bellman optimality. Limited coverage can therefore
yield episode- or trajectory-dependent critic parameters.

\section{Current Precursor: BOCO Critic and Constrained Control Learning}

\textbf{BOCO (Bellman Optimality via Constrained Optimization)} first
parameterizes the unknown optimal critic,
\begin{equation}
  J^*(x)\approx\widehat J_\theta(x),
\end{equation}
and, at each selected update state, expresses the local Bellman minimization
and equality as a finite-dimensional constrained optimization problem. This
makes the local condition amenable to numerical KKT or primal-dual
implementation:
\begin{equation}
  \begin{aligned}
    F_k(u,\theta)
    &:=
    g(x_k,u)
    +
    \alpha\widehat J_\theta
    \!\left(x_{k+1}(u)\right),\\
    \delta_k(u,\theta)
    &:=
    F_k(u,\theta)
    -
    \widehat J_\theta(x_k).
  \end{aligned}
\end{equation}

\begin{equation}
  \begin{aligned}
    \min_{u,\theta}\quad
    &F_k(u,\theta)\\
    \mathrm{s.t.}\quad
    &\delta_k(u,\theta)=0,\\
    &c_i(x_k,u)\le 0,
    \qquad i=1,\ldots,n_c.
  \end{aligned}
\end{equation}

Here $x_{k+1}(u)$ denotes the successor associated with candidate input $u$.
Evaluating this quantity for arbitrary candidate inputs requires a declared
control-usable model, simulator, or other counterfactual mechanism; one
realized transition supplies the successor only for the input that was actually
applied. The set $\mathcal U(x_k)$ contains admissible inputs, $n_c$ is the
number of explicit inequality constraints, and $c_i$ is the $i$th such
constraint. It may represent an actuator limit, a state-dependent feasibility
condition, or a safety-related operating bound. These constraints can be added
explicitly instead of being handled only by a penalty.

A global critic still requires accumulated streaming constraints, a declared
collocation or update set, or another function-approximation condition; one
local equality does not identify the complete value function.

The Lagrangian is
\begin{equation}
  \mathcal L_k(u,\theta,\lambda)
  =
  F_k(u,\theta)
  +
  \lambda_\delta\delta_k(u,\theta)
  +
  \sum_{i=1}^{n_c}\lambda_i c_i(x_k,u),
  \qquad
  \lambda_i\ge 0.
\end{equation}
At online update $t$, the primal control, critic parameter, and multipliers
can be updated conceptually as
\begin{equation}
  \begin{aligned}
    u_k^{\mathrm{BOCO}}
    &\in
    \arg\min_{u\in\mathcal U(x_k)}
    \mathcal L_k(u,\theta_t,\lambda_t),\\
    \theta_{t+1}
    &=
    \theta_t-\eta_{\theta,t}
    \nabla_\theta\mathcal L_k
    \!\left(u_k^{\mathrm{BOCO}},\theta_t,\lambda_t\right),\\
    \lambda_{\delta,t+1}
    &=
    \lambda_{\delta,t}
    +
    \eta_{\delta,t}\,
    \delta_k\!\left(u_k^{\mathrm{BOCO}},\theta_t\right),\\
    \lambda_{i,t+1}
    &=
    \left[
    \lambda_{i,t}
    +
    \eta_{\lambda_i,t}\,
    c_i\!\left(x_k,u_k^{\mathrm{BOCO}}\right)
    \right]_+ .
  \end{aligned}
\end{equation}

Here $\lambda_t=(\lambda_{\delta,t},\lambda_{1,t},\ldots,\lambda_{n_c,t})$
collects the current multipliers, the $\eta$ terms are positive update step
sizes, and $[\cdot]_+$ projects each inequality multiplier onto the
nonnegative orthant. The equality multiplier $\lambda_\delta$ is unrestricted
in sign. The physical state $x_k$ is embedded in $F_k$, $\delta_k$, $c_i$, and
therefore $\mathcal L_k$.

Unlike TD-error reduction alone, BOCO keeps control minimization, critic
consistency, and explicit constraints together. In the current formulation,
$u_k^{\mathrm{BOCO}}$ is the primal actor decision and a separate
policy-network parameter $\phi$ is not required. A future amortized policy
network could be trained to reproduce this solution, but that is a distinct
extension.

The current BOCO paper provides this formulation and constrained nonlinear
simulation evidence. Formal stability, safety, broader task-domain learning,
separate policy-network training, and real-system validation require
additional assumptions and evidence.

\section{Application Domains}

\begin{itemize}
  \item \textbf{Automatic tuning and calibration of control systems for
  electric drives and vehicle or mobility systems}, using operational data to
  improve the deployed critic and controller.
  \item \textbf{Controller learning in complex, unstructured real-world
  environments} where a faithful simulator and exhaustive offline training
  data are difficult to obtain.
\end{itemize}

\section{References}

\begin{itemize}
  \item
  \href{https://kaist-mic-lab.github.io/publications/2025-constrained-optimization/}{
  H. Lee and K. Choi, ``Constrained Optimization Formulation of Bellman
  Optimality Equation for Online Reinforcement Learning,'' International
  Conference on Robot Intelligence Technology and Applications (RiTA), accepted,
  in press, 2025.
  }
  \item
  \href{https://kaist-mic-lab.github.io/publications/2026-online-constrained/}{
  H. Lee and K. Choi, ``Online Constrained Reinforcement Learning for Optimal
  Tracking,'' International Federation of Automatic Control (IFAC) World
  Congress, accepted, in press, 2026.
  }
  \item
  \href{https://kaist-mic-lab.github.io/publications/2025-online-actor/}{
  H. Lee and K. Choi, ``Online Actor Critic Learning for Optimal Tracking in
  Servo Positioning Systems,'' Annual Conference of the IEEE Industrial
  Electronics Society (IECON), pp. 1--7, 2025.
  }
\end{itemize}

% ========================================
%         Appendix
% ========================================

% \begin{appendices}
% \end{appendices}

\addcontentsline{toc}{chapter}{Bibliography}
\bibliographystyle{IEEEtran}
\bibliography{refs}

\end{document}
